\section{Raspberry Pi-integrasjon og unntakshåndtering, 2.11.2026}
\label{sec:piunntak}
\statusbox{\textbf{Verifisert i denne økten:} prosjektet er hentet fra
arbeidsgrenen med grunnrevisjon \kode{cf1178c} og bygget på Raspberry Pi 5.
1857 Ada-kontroller og 35 Python-tester består. Motortransporten er simulert
i de nye feiltestene. Ingen micro:bit var tilkoblet, ingen firmware er flashet,
og ingen ny fysisk kjøring er gjennomført. Den tidligere I2C-oppstartsfeilen
er fortsatt en åpen fysisk verifikasjon.}

\subsection{Avgrensning og arkitektur}
Raspberry Pi brukes som utviklings- og loggmaskin. En USB-datakabel går til
micro:bitens USB-grensesnitt; sanntidsstyringen forblir skrevet i Ada på
micro:bit V2. Dette følger produsentens dokumenterte USB-støtte, som også
omfatter Raspberry Pi~\cite{microbitusb2026}. Motor- og sensorforsyningen
skal kontrolleres før aktiv testing. Pi-ens GPIO kobles ikke til motorene.
Pinneplanen i \kode{PI-OPPKOBLING.md} viderefører siste dokumenterte oppsett;
den er ikke en ny fysisk bekreftelse.

Applikasjonen er fortsatt sekvensiell. Sensor-, navigasjons- og motorpakker
samarbeider gjennom én styringsflyt; stoppknappen setter et atomisk flagg.
En fremtidig taskarkitektur må ha eksplisitt eierskap til ultralyd og I2C,
tidsstemplede data og beskyttet kommunikasjon. Det nye feilobjektet har én
sekvensiell eier og skal ikke brukes samtidig fra flere tasks uten omarbeiding.

\subsection{Exceptions som avgrenser faktiske feil}
Ada overfører kontrollen til en passende handler når en exception oppstår;
en ubehandlet exception i en task behandles annerledes enn propagation i en
vanlig subprogramkjede~\cite{adaexceptions2026}. Prosjektet bruker runtime
\kode{embedded-nrf52833}. Støtte må vurderes mot dette konkrete valget;
lettere runtimes kan begrense propagation~\cite{adart2026}.
\kode{Adaptive_Errors} skiller motorbussfeil, ufullstendig initialisering,
servogrenser, latchet stopp og feil ved stoppknappens registrering. Innebygde
\kode{Constraint_Error}, \kode{Program_Error} og \kode{Storage_Error}
klassifiseres separat. En vanlig hindring eller en ugyldig avstand behandles
fortsatt som sensordata/FSM-hendelse, ikke som en exception per avlesning.

Ved I2C-feil lagres stoppforespørselen og videre kjøring avvises.
Feilmeldingen inneholder status, enhetsadresse og register. Motorstyringen
stopper motorutgangene \emph{før} oscillatorens oppvåkning, slik at gamle
PWM-verdier ikke vekkes før stopp. Ved ukjent auto-increment skrives
\kode{FULL_OFF} separat til de åtte motor-PWM-kanalene; dette unngår å anta
at en ufullført oppstart har aktivert burstskriving. Registervalget bygger
på PCA9685-databladet~\cite{pca9685_2026}. Endringen er verifisert med
simulert I2C, men må fortsatt verifiseres på den fysiske motorstyringen.

\kode{Adaptive_Fault_Handler} setter først stoppflagget, forsøker så én
avgrenset stoppoperasjon og skriver deretter diagnostikk. En sekundær
stoppfeil bevares ved siden av den opprinnelige feilårsaken. En loggerfeil
kan ikke utløse ny kjøring eller gjentatte gjenopprettingsforsøk. Det finnes
ingen automatisk reset av feiltilstanden. Hovedprogrammet og diagnostikk V2
bruker samme handler. Dette erstatter tidligere skjult \kode{when others}
ved feil i selve stoppforsøket. Feilgrensen dekker exceptions i den
eksekverende styringsflyten; den erstatter ikke en last-chance-handler ved
feil under bibliotekelaborering eller ved maskinvarelås.

\begin{table}[htbp]
\centering\small
\caption{Feilrespons og begrensninger i den nye implementasjonen.}
\begin{tabularx}{\textwidth}{@{}P{34mm}YY@{}}\toprule
Hendelse & Implementert respons & Hva responsen ikke beviser\\\midrule
Motor-I2C avviser sending & Navngitt exception og stoppflagg & Årsak til elektrisk bussfeil\\
Ukjent oppstartstilstand & Enkeltvis FULL\_OFF før oppvåkning & Faktisk PWM eller stillstand på bilen\\
Stopp avvises av bussen & STOP UNCONFIRMED og bevart sekundærfeil & At programvaren kan stoppe uten fungerende buss\\
Knapp under motorsending & Ny kontroll etter sending, så stopp & Elektrisk uavhengig nødstopp\\
Ugyldig servovinkel & Ingen sending, exception til supervisor & At alle tillatte vinkler er mekanisk klarert\\
\bottomrule\end{tabularx}
\end{table}

\subsection{Tidsmessige konsekvenser}
Motor og OLED deler ekstern I2C og unngår nå unødvendig ny initialisering av
samme buss. OLED-skriving har eksplisitt 20 ms timeoutparameter, som motoren.
Normal motorstopp bruker én bursttransaksjon når auto-increment er kjent.
Etter usikker oppstart kan stopp kreve åtte enkelttransaksjoner:
\[
 B_{\mathrm{stopp,IO}} \leq 8\cdot 20\,\mathrm{ms}=160\,\mathrm{ms}.
\]
Dette er et betinget transaksjonsbudsjett i den patchede driveren. Scheduler,
programinstruksjoner, sensorventing og fysisk bremsing kommer i tillegg.
Ved første avviste transaksjon avbrytes forsøket og stopp er ubekreftet.
Det er derfor ikke dokumentert en 160 ms end-to-end-deadline.
Knapp B gir en programvareforespørsel; en ytre strømbryter må fortsatt være
tilgjengelig ved fysiske forsøk. Fremtidig RTA må bruke den faktiske
arkitekturen og omfatte blokkering og eventuelle self-suspensions.

\subsection{Testresultat og faglig sporbarhet}
De eksisterende 1753 Ada-kontrollene består fortsatt. De 104 nye kontrollene
omfatter klassifisering, rekkefølgen stoppflagg--stopp--logging, sekundær
stoppfeil, loggerfeil, alle tolv skrivepunkter i oppstarten, servogrenser,
stopp under motorsending og avvisning av ny kjøring etter feil.
Testene kjører faktisk applikasjonskode og DFR0548s PWM-koding mot en
simulert I2C-transport. De utgjør ikke 104 fysiske forsøk eller et formelt
sikkerhetsbevis. ARM-bygg av hovedprogram, V2 og GPIO V3 består.

Det passive USB-verktøyet \kode{pi_link.py} identifiserer USB-enheten og
samler tidsavgrensede rålogger uten å sende kommandoer. Mottakstid på Pi-en
er ikke micro:bitens utførelsestid. Ingen firmwareidentitet eller fysisk
motorstatus utledes av en åpen port. De 35 Python-testene består av 29
eksisterende tester og seks nye tester av deteksjon, avgrensning, fravær av
skrivekommandoer og bevaring av delvise logger.

Oppgaveteksten og eksisterende rapport brukes som faglig arbeidsgrunnlag.
Selve forelesnings-PDF-en og gjeldende vurderingskriterier er ikke tilgjengelig
i dette checkoutet. Detaljert pensumdekning avventer disse kildene.
\kode{PENSUM-OG-DEMOPLAN.md} kobler kode, teori og evidens og angir videre
målinger. Hovedrapporten bør ved sluttredigering tilpasses oppgavens
anbefalte 15--20 sider, mens utviklingshistorikk og vesentlig kode legges i
vedlegg. Ingen påstand om sensorvurdering eller karakter kan utledes av
antall funksjoner eller tester alene.
